\subsubsection{SPV Bribe Attacks: Derivations \& Breakpoints}


Let's assume that an attacker is willing to bribe miners to exclude a fraud proof.
The goal of the attacker is to continue this attack beyond some breakpoint, so how much will that cost them over the course of the attack?

First, let's establish a baseline: what does a healthy simplex look like?
\emph{Healthy} meaning: no fraud proof is possible because all blocks and state-transitions are valid.
We need to know what a healthy simplex (and its PoR graph) looks like because a successful attack will be indistinguishable.

The defining quality of a healthy simplex is that additions to the PoR graph are highly interconnected.
When new blocks are mined, they reflect all possible remote blocks and link back to all possible local blocks.
Thus, there should be no partitions or asymmetric boundaries.

\mk{"The asymmetric boundary lies" -- huh? check for making sense. boundary is inside the graph }

An \emph{asymmetric boundary} emerges when a group of miners contributes \emph{exclusively} to a \emph{subgraph} of the full PoR graph (e.g., a subgraph that censors the fraud proof).
The asymmetric boundary lies between the full (honest) PoR graph and the dishonest subgraph.



When an invalid block is mined, individual nodes of the network should rapidly become aware of the fraud proof --- and we expect this if the P2P layer is working correctly.
However, during an ongoing (successful) attack, that fraud proof is not recorded in the PoR graph.
The fraud is `invisible' when considering just the blockchain data itself, which means that SPV proofs may be fraudulent.

When fraud is detected, the PoR graph should proceed through multiple phase changes until normality is restored (\autoref{fig:por-graph-phases}).

\begin{figure}[h]
    \begin{center}
        \tikzstyle{n} = [align=center, draw, rectangle, rounded corners, inner sep=5pt]
        \tikzstyle{ar} = [->, anchor=north, rounded corners, shorten >= 2pt, shorten <= 2pt]
        \tikzstyle{label} = [align=center, anchor=center]
        \begin{tikzpicture}[node distance=1.8cm]
            % states of PoR graph
            \node (s0) [n] {0. PoR Graph \\ Looks Healthy};
            \node (s1) [n, right=of s0, xshift=.95cm] {1. Attack \\ In Progress};
            \node (s2) [n, right=of s1, xshift=.6cm] {2. Partial FP \\ Acknowledgment};
            \node (s3) [n, below=of s1.north, yshift=.2cm] {3. Complete FP \\ Acknowledgment};

            % group s0 and s1
            \node[n, fit=(s0) (s1), dashed, inner sep=7pt] {};

            % arrows and labels
            \draw [ar] (s0) -- node [label] {Fraudulent \\ block mined} (s1);
            \draw [ar] (s1) -- node [label] {FP first \\ recorded} (s2);
            \draw [ar] (s2) -- (s2 |- s3) -- node [label] {FP `pollutes' \\ PoR graph} (s3);
            \draw [ar] (s3) -- node[label] {PoR graph now \\ resynchronized} (s3 -| s0) -- (s0);
        \end{tikzpicture}
    \end{center}
    \caption[
        Phase transitions of a PoR graph when fraud occurs.
    ]{
        The phase transitions of a PoR graph when fraud occurs.
        The dashed outline around phase 0 and 1 indicates that it's not possible to observe this phase change on-chain.
        Nodes connected to the P2P network should be aware of the fraud, but no chain has yet included the fraud proof in the PoR graph.
    }
    \label{fig:por-graph-phases}
\end{figure}

An attacker always fails if the network reaches phase 3.
The attacker's primary goal is to prevent the network transitioning from phase 1 to phase 2.

The dynamics of the attack change for both the attacker and each miner during each phase.

For example, there is a qualitative difference between phase 1 and phase 2, and it is in the attacker's interest to prevent the network transitioning to phase 2 if possible.
Therefore it is most valuable for the attacker if they can prevent the FP being recorded.
It is also more profitable for miners, since the size of bribes grow quadratically.

Once the network has reached phase 2, it is only a matter of time before the network reaches phase 3, at which point we are back at the start.

If the network is kept in phase 1, then the PoR graph and (nearly all) chain data is indistinguishable from a healthy simplex.



\subsubsubsection{Phase 1 $\not\to$ Phase 2 on the Target Chain}

The first honest block on the target chain must, by definition, link back to the fraudulent block as an invalid ancestor.
Let’s assume that there’s a reward $(C_r H)$ for publishing the fraud proof, where $C_r$ is the
block reward of the target chain, and $H$ is the proportion of the block reward each miner gets.
The additional reward for this block after $k$ attacking blocks is $C_r H k$.
Thus, after $H^{-1}$ blocks, the \emph{effective} block reward for an honest miner will have doubled.

A new honest block thus becomes heavily incentivized over time.

To mitigate this, the attacker must raise the difficulty on the target chain such that all other potential miners find it comparatively unprofitable.

Compared to mining honestly, the attacker forgoes the reward of $C_r H k$ for the k-th block.
That reward is lost to excessive work securing the target chain which would otherwise earn the attacker that much revenue from publishing the fraud proof.
Additionally, the attacker loses the block rewards each time --- so these losses are $C_r (1 + H k)$ per block after $k$ blocks.
Let $C_r \cdot l_m$ represent the total losses due to mining, and $z$ the number of blocks for the attack: factoring out $l_m$ in this way allows us to analyze the situation independent of the block reward.
In order for the losses $l_m$ to be greater than the rewards, we must have:
\begin{align}
    l_m &> \sum\limits_{k=0}^z \Bigl( 1 + H k \Bigr) \nonumber \\
    &= z + 1 + H \sum\limits_{k=0}^z k \nonumber \\
    &= z + 1 + H \frac{z (z + 1)}{2}
    \label{eq:spv-atk-losses-mining}
\end{align}

Additionally, the rise in difficulty of these blocks presents a practical limit for the attacker.
When $k > q H^{-1} N_1$ the attacker is, \emph{by definition}, unable to keep up the attack.
That is because, at that point, the attacker's losses are $C_r (1 + H (q H^{-1} N_1)) = C_r (1 + q N_1)$.
Assuming an even hash-rate distribution, $C_r N_1 (p + q)$ represents the \emph{sum} of \emph{all} block rewards over the network.
Therefore, the maximum output (measured in coins) of $q$ proportion of the hash-rate, \emph{by definition}, is $C_r N_1 q$.
An attacker must have additional support to maintain this portion of the attack --- when the mining losses per block (approx $C_r H k$) exceed the attacker's mining capacity ($C_r N_1 q$).
% We have:
\begin{align}
    \text{\hspace{8em}}
    && C_r H z &< C_r N_1 q \nonumber \\
    && z &< q H^{-1} N_1 & H \ge {N_1}^{-1} \nonumber \\
    && z &< q {N_1}^2 & \text{(worst case)} \nonumber
\end{align}
If $N_1 = 100$, then no attack is viable (without help) after 5000 blocks, unless $q > 0.5$.
At that point, at least 50\% of the network hash-rate needs to be dedicated to the target chain to make it unprofitable enough to dissuade honest miners.

\subsubsubsection{Phase 1 $\not\to$ Phase 2 and $\BribeCost(z)$}
\label{sec:bribe-stop-1-to-2}

To prevent the network transitioning from phase 1 to phase 2, the attacker incurs a cost.
We can define this cost \emph{in terms of the block reward}.
Let the total cost to the attacker, in the $k$-th round, be $C_r Z_k$ coins.
($Z_k$ is the cost in block rewards.)

\begin{itemize}
    \item Each miner (on each simplex-chain) needs bribing (but only if they make a block).
    \item Each round, the attacker must pay $C_r Z_k$ in bribes (where $Z_k$ is the bribe as a multiple of the block reward, and $k$ is the round, starting at 1).
    \item These bribes are given to $(N_1 - 1)$ miners (we subtract 1 to exclude the target chain from this split). One bribe per R-chain.
    \item However, an attacker with $q$ proportion of the hash-rate can recover $q(N_1 - 1)$ of these bribes, at most, and thus loses at least $p(N_1 - 1)$ of them.
    \item A miner receives $C_r H k$ coins as a reward for publishing a fraud proof that invalidates $k$ blocks.
    \item Each round, miners can choose to publish the fraud proof, which earns them $C_r H k$ in rewards.
    \item Thus, each round, miners have an opportunity cost of $C_r H k$ if they do \emph{not} take the bribe.
    \item This only makes rational sense if the miner's bribe is greater than their opportunity cost: $C_r Z_k / (N_1 - 1) > C_r H k$; \emph{and} the miner is \emph{guaranteed} to receive the bribe, regardless of whether the attack succeeds.
    \item The attack lasts $z$ rounds.
    \item Let $C_r \cdot l_b$ be the total losses due to bribes.
\end{itemize}

Therefore, for the bribes of round $k$ to be more than the rewards on all chains, we must have:
\begin{align}
    % \text{let } H(z) &= \frac{z}{N_1} \nonumber \\
    % \implies N_1 H(z) &= z \nonumber \\
    % Z_1 &> H (N_1-1) \nonumber \\
    % Z_2 &> 2 H (N_1-1) \nonumber \\
    Z_k &> k H (N_1-1) \nonumber \\
\end{align}
Let $l_b = p (Z_1 + \cdots + Z_z)$ denote the expected cost of the attack for the $z$ rounds. Substituting in the above:
\begin{align}
    l_b &> p H (N_1-1) \sum\limits_{k=1}^z k \nonumber \\
    l_b &> p H (N_1-1) \frac{z (z+1)}{2}
        \label{eq:spv-atk-losses-bribes}
\end{align}

% \begin{align*}
%     (q + N_1 p) H \frac{z(z+1)}{2}
% \end{align*}

Let $\BribeCost(z)$ denote the minimum cost of the attack for $z$ rounds, given our current understanding.
Combining \autoref{eq:spv-atk-losses-mining} with \autoref{eq:spv-atk-losses-bribes}, and noting that $p+q=1$ yields:
\begin{equation}
    \BribeCost(z)= l_b + (p+q)l_m > H (N_1 p + q) \frac{z(z+1)}{2} + z + 1
    \label{eq:spv-atk-losses-total}
\end{equation}

% \todo{nomenclature clash: $h$ was previously used both for $H(Np+q)$ and in $h+d=1$. set $\eta = H(Np+q)$ as a replacement. keep $h+d=1$.}

Observe that this is a lower bound on $\BribeCost(z)$.
To reduce complexity, we'll also define \BribeCost \emph{in terms of the block reward} (i.e., in units of \emph{block rewards} rather than \emph{coins}).
We can come back to it later if we need to.
\begin{align}
    \BribeCost(z) &= \eta \frac{z(z+1)}{2} + z + 1
        & & \text{where } \eta = H(N_1 p + q)
    \label{eq:spv-bribe-cost-1}
\end{align}
We should also note the inverse of \BribeCost, which gives us a lower bound on the number of blocks for which we should wait, given a total outgoing value, $v$ (expressed as a multiple of the block reward).
\begin{align}
    \BribeCost^{-1}(v) &= \frac{\sqrt{
        (\eta - 2)^2 + 8 \eta v
    } - \eta - 2}{2\eta} & & \text{where } \eta = H(N_1 p + q)
    \label{eq:spv-bribe-cost-inverse}
\end{align}

\begin{comment}

\subsubsubsection{Domain and Range of $\BribeCost(z)$ in the SRT Context}

It is clear that $\BribeCost(0) = 1$.
That is: the attacker has lost exactly 1 block reward (the reward for the fraudulent block).

Since this is a simplistic model, we'll only consider non-negative integer values of $z$.

Although it is clear that $\lim_{z \to \infty} \BribeCost(z) = \infty$, there is a practical upper limit.

Recall that we introduced \BribeCost via the requirement that $v - \BribeCost(z) < 0$ assuming an SPV attack yields $v - \BribeCost(z)$ profit for the attacker.

There is a practical upper limit to $v$: \emph{all} of the root tokens on that simplex-chain.

In the \emph{Single Root Token} (SRT) PoR conversion context, recall that we defined the block reward in \autoref{eq:srt-reward} as:
\begin{align*}
    C_r &= \frac{C_t \cdot I}{G_t \cdot C_f} & & \frac{\text{C-coins}}{\text{C-block}}
\end{align*}
Currently, $I$ (the global inflation rate) is measured in coins/s.
Let's assume that the inflation rate is proportional to $G_t$ (the total supply of root tokens).
For example, we might desire an $I$ such that the \emph{annual} inflation amount is 2\% of $G_t$.
Let's use $f$ to refer to that percentage of $G_t$ (e.g., $f = 0.02$), and $\syr = 365.25 \times 86400 = 3.078 \times 10^{7}$ seconds in a year.
Thus, we have:
\begin{align}
    I &= \frac{f G_t}{\syr}
        & & \FracUnits{coins}{second} \nonumber \\
    \implies C_r &= \frac{f G_t}{\syr} \cdot \frac{C_t}{G_t \cdot C_f}
        & & \FracUnits{C-coins}{C-block} \nonumber \\
    C_r &= \frac{f C_t}{\syr C_f} \label{eq:block-reward-to-ct}
\end{align}
To convert $v$ from coins to block rewards, we divide by $C_r$.
The maximum $v$ possible is $C_t$.
\begin{align*}
    \frac{v}{C_r} &= \frac{v \syr C_f}{f C_t} \\
    \text{Set } v = C_t \implies \frac{C_t}{C_r} &= \frac{\syr C_f}{f} & \text{which is a constant.}
\end{align*}
So, the upper limit of $v$, expressed as a multiple of block rewards, is a constant.
Particularly, it is the number of blocks in a year divided by the inflation rate $f$.
This represents the upper limit of the range of $\BribeCost(z)$ when using the SRT context.

Let's apply \autoref{eq:spv-bribe-cost-inverse} to get an idea of what the maximum waiting time might be.
For variables, let's use $q = p = 0.5$, $N_1 = 100$, and $H = 1/N_1 = 1/100$, which implies $\eta = 0.505$; and $C_f = \nicefrac{1}{15}$ and $f = 0.02$, which implies $\syr C_f f^{-1} = 1.026 \times 10^8$.
\begin{equation*}
    \BribeCost^{-1}(\syr C_f f^{-1}) \approx 20 115.3 \text{ blocks}
        \implies \sim 3.5 \text{ days}
\end{equation*}
If, instead, $H = 1/\sqrt{N_1}$, then $\eta = 5.05$ and
\begin{equation*}
    \BribeCost^{-1}(\syr C_f f^{-1}) \approx 6373.8 \text{ blocks}
        \implies \sim 1.1 \text{ days}
\end{equation*}

\end{comment}

\subsubsubsection{Phase 1 $\to$ Phase 2}

The network transitions from phase 1 to phase 2 as soon as a first miner produces a block that includes the fraud proof.

Once this happens, an asymmetry emerges: honest blocks (which acknowledge the fraud proof) produce proofs of reflection that are \emph{not} usable by the dishonest miners.
If the dishonest miners tried to use those PoRs, they'd produce invalid blocks unless they also acknowledge the fraud proof.
Thus, using those PoRs effectively turns dishonest miners into honest miners.



\subsubsubsection{Phase 2 $\to$ Phase 3}

The duration of this transition depends on the proportion of honest miners: $h$.
Let's define $h$ via $h + d = 1$, where $d$ is the proportion of dishonest miners, including the attacker.\footnote{
    Dishonest miners are those which will accept a bribe and participate in the fraud proof censorship attack, but are not necessarily controlled by the attacker directly.
    We assume that there are some miners that would accept a bribe for \emph{this} attack, but not for \emph{any} attack.
}
Since $d$ includes the attacker, we can say $d > q$, and $h < p$.

Similar to \autoref{sec:dos-and-dags}, the asymmetry between honest and dishonest miners means that, even for a small $h$, the chain-weight of the honest chain-tips will eventually dominate the chain-weight of dishonest chain-tips.
After that point, each simplex-chain's main chain will include the fraud proof, and the dishonest chain-tips will never be able to out-weigh the honest ones.

Let us use a vector to represent the weight of the honest and dishonest chain segments (relative to their most recent common ancestor).
This is purely for convenience so that we can analyze both at the same time.
The top row is the weight of the honest chain segment, and the second row is the weight of the dishonest chain segment.
At the start of the attack ($t_0$), neither chain segment has any blocks, so $t_0 = \vec{0}$.

Each round (block period), we expect that $h$ proportion of new blocks are honest, and $d$ proportion are dishonest.
So, after 1 block period, we expect:
\begin{equation*}
    t_1 = t_0 + \vecXY{N_1 h}{N_1 d} = N_1 \vecXY{h}{d}
\end{equation*}

After 2 block periods, the $h$ proportion of new blocks that are honest can use PoRs via the dishonest blocks.
Due to the asymmetry, the dishonest blocks can not.
\begin{equation*}
    t_2 = t_1 + \vecXY{N_1 (h+d)}{N_1 d} = N_1 \vecXY{h}{d} + N_1 \vecXY{1}{d} = N_1 \vecXY{(h + 1)}{2 d}
\end{equation*}

\mk{change $k$ below to now overload general meaning}
Each round, the honest chain segments grow by the weight of $N_1$ blocks, but the dishonest chain segments grow by $N_1 d$ blocks.
Therefore, at the $k$-th round, $k \ge 2$:
\begin{align}
    t_k &= t_{k-1} + N_1 \vecXY{1}{d} = t_{k-2} + 2 N_1 \vecXY{1}{d} \nonumber \\
    \text{observe: } t_k &= t_{k-i} + i N_1 \vecXY{1}{d} & \text{for some } i \in [1, k) \nonumber \\
    \text{set } i = k - 1 \implies
    t_k &= t_1 + N_1 (k-1) \vecXY{1}{d} \nonumber \\
    t_k &= N_1 \vecXY{h}{d} + N_1 (k-1) \vecXY{1}{d} \nonumber \\
    \therefore t_k &= N_1 \vecXY{h + k - 1}{dk} \label{eq:h-to-d-chain-segs}
\end{align}

Thus, the honest chain-tips dominate when
\begin{align*}
    h + k - 1 &> dk \\
    k - dk &> 1 - h \\
    k(1 - d) &> 1-h \\
    h k &> 1-h \\
    k &> h^{-1} - 1
\end{align*}

If $h = \nicefrac{1}{N_1}$, i.e., for each honest miner there are $(N_1 - 1)$ dishonest miners, then this inequality becomes: $k > N_1 - 1$.
% \begin{equation*}
%     k > N_1 - 1
% \end{equation*}

If only 1\% of miners refuse the bribe ($h = 0.01$), and that the other 99\% of miners continue to mine only blocks that do not include the fraud proof, then we have: $k > 99$.

If we assume an honest minority of substantial size, e.g., $h = \nicefrac{1}{3}$, then $k > 2$.
